\scanpage{120}
% archive sheet number: 119

En pratique, il sera vrai que $\handwritten{\overline}{U'}$ est schématiquement dense dans $\handwritten{\overline}{X'}$\inserted{et $\handwritten{\overline}{Y'}$ séparé sur $\handwritten{\bar}{Y}$}, de sorte que $\handwritten{\overline}{f'}$ sera déterminé de façon unique en termes des autres données, comme l'unique prolongement de $\handwritten{\overline}{U'}\handwritten{\longrightarrow}\handwritten{\overline}{V'}$ induit par $f'$; ce sera le cas en particulier dans le cas envisagé plus haut où $\handwritten{\overline}{X'}$ se déduit de \deleted{$\handwritten{\bar}{\unsure{U}}$} $\handwritten{\overline}{U'}$ par normalisation. Bien entendu\deleted{\unsure{r}}, pour pouvoir conclure\inserted{\handwritten{alors}} que tout sous-groupe d'inertie de $G$ rel. à $x$ s'envoie dans un sous-groupe d'inertie de $H$ rel. à $y$, il faudra avoir vérifié l'\uline{existence} d'un p\deleted{p}rolongement $\handwritten{\overline}{f'}$ de $\handwritten{\overline}{U'}\handwritten{\longrightarrow}\handwritten{\overline}{V'}$\inserted{\handwritten{(la commutation aux actions de $G$ dev. alors \uline{automatique})}}. \handwritten{cela \unsure{marchera}, en interprétant $\unsure{\overline{f'}|f'}$ comme $\overline{X'}\unsure{\longrightarrow|-}\overline{Y'}\mathbin{\times}_{\bar Y}\overline{X'}$, lorsque $\overline{Y'}$ est \uline{entier} \deleted{\illegible{1}} sur $\bar Y$, cf. EGA II 6.1.14.}

\handwritten{3.} Revenons à la situation de 1), supposons que $\underline{O}_{X,x}$ soit un anneau local régulier, et que $\handwritten{\bar}{U}$ soit le complémentaire dans le schéma stricte\deleted{\unsure{men}}ment local régulier $\handwritten{\bar}{X}$ d'un diviseur \inserted{$\handwritten{\bar}{D}$} à croisements normaux. On est en particulier dans le cas favorable où la définition de $\handwritten{\overline}{X'}$ est immédiate par normalisation. Supposons que la ramification de $\handwritten{\overline}{U'}$ en les points maximaux de $\handwritten{\bar}{D}$ soit \uline{modérée}. On sait alors, par le calcul local bien connu, que pour tout $\handwritten{\unsure{\xi|\zeta}\in}\handwritten{\overline}{X'}_{\handwritten{\bar}{x}}$, le groupe d'inertie correspondant $G$ s'identifie canoniquement \deleted{au groupe} à un groupe quotient du groupe
\margin{\handwritten{$C=$ l'ensemble des composantes irred. de $\bar D$}}
\[
I=\handwritten{\prod}_{\handwritten{\unsure{\ell|l}\ne}p}\underline{Z}_{\handwritten{\ell}}(1)(k(\handwritten{\bar}{x}))^C=\handwritten{\prod}_{\handwritten{\ell\ne}p}T_{\handwritten{\ell}}(k(\handwritten{\bar}{x}\unsure{)}\ )^C\ ,
\tag*{(3.1)}
\]
$p$ étant la caractéristique résiduelle en $x$. De p\deleted{\unsure{m}}lus, ce même calc\uline{u}l local montre que le facteur de ce groupe d'inertie correspondant à la composante $\handwritten{\bar}{D}_c$ $(c\handwritten{\in}C)$ du diviseur de ramification $\handwritten{\bar}{D}$ coïncide aussi avec le groupe d'inertie relatif à n'importe quel point de $\handwritten{\overline}{X'}$ au dessus du point maximal de $\handwritten{\bar}{D}_c$ (ou de tout autre point de $\handwritten{\bar}{D}_c$ qui n'appartient à aucune autre composante irréductible de $\handwritten{\bar}{D}$) qui générise $\handwritten{\xi}$ ; comme il y a toujours de tels points, cela montre que la classe de conjugaison du facteur en question peut aussi s'interpréter comme la \deleted{grou} classe des groupes d'inertie en les points envisagés de $\handwritten{\bar}{D}_c$ ; en particulier, cette classe de conjugaison ne dépend que de la composante irréductible de $D=X-U$ \deleted{, image} de point générique l'image du point générique de $\handwritten{\bar}{D}_{\handwritten{c}}$.
